\documentclass[11pt]{article}
\usepackage[T1]{fontenc}
\usepackage[utf8]{inputenc}
\usepackage{amsmath, amssymb}
\usepackage{booktabs}
\usepackage{geometry}
\geometry{margin=2.6cm}
\usepackage{setspace}
\onehalfspacing
\usepackage[bottom]{footmisc}
\usepackage[hidelinks]{hyperref}
\usepackage{threeparttable}
\usepackage{float}
\usepackage{titlesec}
\titleformat{\section}{\normalfont\large\bfseries}{\thesection.}{0.6em}{}
\setlength{\parskip}{2pt}

\title{\textbf{Do Environmental Provisions in Trade Agreements\\Reshape Export Composition?}\\[8pt]
\large Extended summary}
\author{Edoardo Vitella\\\small University of Trento \& Free University of Bozen}
\date{September 2026}

\begin{document}
\maketitle
\vspace{-18pt}

% ---------------------------------------------------------------------------
\section{The question}

Nearly every trade agreement signed in the past two decades contains environmental language, and
the number of such provisions has grown steadily. Whether any of it changes what countries actually
trade is close to an open question. The evidence we have comes almost entirely from gravity models
on aggregate bilateral flows, which can tell us whether trade between two countries grew, but not
whether its composition shifted --- whether a country began exporting more environmental goods and
fewer pollution-intensive ones.

This paper asks that second question for the world's largest exporter. Between 2000 and 2015 China
signed fourteen preferential trade agreements covering twenty-five destination economies, with
environmental content ranging from a handful of generic cooperation clauses to a dedicated chapter.
Did the agreements with more environmental content tilt Chinese exports toward green goods and away
from dirty ones?

The answer is no on the green margin, with a bound attached, and inconclusive on the dirty margin.
The more useful contribution is not the null itself but what produces it: the environmental content
of these particular agreements consists almost entirely of language with no mechanism through which
it could move a product mix. The null is what the existing literature's own heterogeneity results
predict for a sample composed this way.

% ---------------------------------------------------------------------------
\section{Setting and data}

\textbf{Trade data.} The universe of Chinese customs transactions, 2000--2015: 45.8 million
firm--HS6--destination--year observations covering 462,651 exporting firms. The granularity is what
makes the design possible --- the same firm can be observed shipping the same product to a
destination covered by an agreement and to one that is not, in the same year.

\textbf{Environmental content.} Two independent codings, used in parallel throughout. The World Bank
Deep Trade Agreements database gives a count of environmental provisions in force with each
destination at each date, together with counts for the non-environmental policy areas. The TREND
database codes environmental provisions at finer granularity, with sub-indices for green market
access, enforcement, hard and soft obligations, and regulatory space. Using two codings guards
against the idiosyncrasies of either.

\textbf{Product classification.} Green goods follow the OECD Combined List of Environmental Goods;
dirty goods follow the standard pollution-intensity taxonomy mapping polluting industries to HS6
codes. Everything else is neutral and serves as the within-cell comparison group.

\textbf{Samples.} Hong Kong and Macao are excluded from the main estimates. Both have Closer Economic
Partnership Arrangements with the mainland, and both are entrep\^ots: a large share of what the data
record as Chinese exports to them is transshipped onward, so the recorded destination and the actual
final market differ systematically. They account for 24.4 percent of treated observations and 50.1
percent of treated export value, and results including them are reported separately. After iterative
removal of singleton observations the estimation sample is 21,519,511 observations with 225
destination clusters, 23 of them treated.

% ---------------------------------------------------------------------------
\section{Empirical strategy}

\subsection*{Why composition rather than levels}

The obvious specification --- regress log exports on environmental depth --- cannot work here, for a
structural reason rather than an empirical one. Any fixed-effect structure that takes bilateral
confounding seriously requires a destination--year intercept, and that intercept absorbs everything
that varies at the destination--year level, environmental depth included. With thirteen distinct
agreement profiles and almost no within-destination change over time, ``having deep environmental
content'' and ``having a deep agreement'' are empirically the same variable.

An interaction survives the destination--year intercept because it varies within the destination--year
cell, across products. That is why the design is built on composition: it compares how green and dirty
products respond relative to neutral ones, inside the same firm--destination--year.

\subsection*{The specification}

\begin{equation*}
\ln x_{fpdt} = \beta_1\, EP_{dt}\!\times\!g_p + \beta_2\, EP_{dt}\!\times\!b_p
+ \gamma_1\, TD_{dt}\!\times\!g_p + \gamma_2\, TD_{dt}\!\times\!b_p
+ \theta_{fpd} + \theta_{fdt} + \theta_{pt} + \varepsilon_{fpdt}
\end{equation*}

$g_p$ and $b_p$ are the green and dirty indicators, $EP_{dt}$ environmental depth, $TD_{dt}$
non-environmental agreement depth. The three sets of absorbed effects do different work.
$\theta_{fpd}$ removes time-invariant features of each firm--product--market relationship.
$\theta_{fdt}$ is what drives identification: it removes everything time-varying and bilateral ---
the agreement itself, demand shocks, exchange rates, the non-random selection of partners --- leaving
only the differential response of green and dirty products within the cell. $\theta_{pt}$ removes
global product-level shocks such as the worldwide expansion of the solar market.

Because environmental depth and general agreement depth are highly correlated, every specification
controls separately for non-environmental depth, and the paper reports every result under two
independent measures of it: the World Bank count and the DESTA index. This turns out to matter, and
Section~5 explains why.

\subsection*{Inference}

Twenty-three treated destinations, sharing thirteen distinct agreement profiles, because the ten
ASEAN partners enter through a single agreement. Under those conditions cluster-robust standard
errors over-reject, sometimes badly. Every estimate is therefore reported with a wild cluster
bootstrap $p$-value ($B = 9{,}999$), which is the reference inference, and the collapsed-panel
estimates additionally with a permutation $p$-value.

The permutation test is worth describing because it does something the bootstrap does not. The
environmental profile of each agreement --- its depth and the years in which it takes effect --- is
detached from its destination and reassigned at random among the treated destinations; the model is
re-estimated; the exercise is repeated a thousand times. The resulting $p$-value is the share of
random reassignments producing a coefficient at least as large as the observed one. It answers the
question: would this number have appeared if we had simply relabelled which of these countries signed
environmental chapters?

% ---------------------------------------------------------------------------
\section{Results}

\subsection*{First, the fixed effects earn their keep}

Estimating the same interactions under progressively richer fixed-effect structures produces a
striking pattern. Under any structure that lacks an absorber with both a destination and a time
dimension, the interactions are frequently positive and sometimes significant: the World Bank dirty
interaction reaches $+0.0064$ ($p = 0.019$), the TREND green interaction $+0.0021$ ($p = 0.021$,
column~3).
Both margins moving in the same direction is not a composition effect --- it is the growth of Chinese
exports to partner countries around the entry into force of an agreement, loading onto whichever
interaction happens to be correlated with agreement timing.

Adding firm--destination--year effects removes it. The World Bank dirty coefficient goes from
$+0.0064$ to $-0.0044$ and the TREND green coefficient from $+0.0018$ to $-0.0001$, with the
significance disappearing alongside the sign. We cannot identify which destination-specific shock was
responsible, only that it had a destination--time structure, since absorbing that structure is what
eliminates it.

\subsection*{The green margin: a bounded null}

On the full panel the green interaction is $-0.0023$ (s.e.\ 0.0039) with the World Bank index and
$-0.0001$ (0.0010) with TREND. Changing the depth control moves the third decimal. The joint $F$-test
on all four composition interactions gives $p = 0.31$ and $p = 0.71$. Quantity and unit value are
equally silent, which matters more than it appears: it rules out the possibility that a volume
response and a price response are cancelling inside the value aggregate. Neither channel moves.

The null holds when the comparison group is replaced three separate times --- comparing products only
within the same HS4 heading, matching destinations on pre-treatment observables, and dropping
untreated destinations altogether so that identification comes from differences in depth among
partners. Across the twelve green cells of the three subsamples the coefficient stays between $-0.0023$ and
$+0.0005$ and its bootstrap $p$-value never falls below 0.58; including the two reference rows, the
range widens only to $[-0.0046, +0.0018]$ and the lowest $p$-value is 0.39.

The bound is what makes this more than a failure to reject. The 95 percent bootstrap interval runs
from $-0.0353$ to $+0.0355$ per World Bank provision. Any effect larger than about three percent per
provision is ruled out at 95 percent confidence; anything smaller, the design cannot see. The
statement the paper defends is a bound, not an absence.

\subsection*{The dirty margin: negative everywhere, decisive nowhere}

The dirty coefficient is negative in almost every specification and its magnitude is stable. Its
statistical standing is not, and the reason is instructive.

Under the World Bank measure of general agreement depth the full-panel coefficient is $-0.0044$ with
an asymptotic $p$ of 0.052, which the bootstrap raises to 0.19. Under the DESTA measure the same
coefficient is $-0.0056$ with asymptotic $p$ of 0.039 and bootstrap $p$ of 0.049. Same data, same
fixed effects, same sample. What differs is collinearity: net of the fixed effects, environmental
depth correlates 0.959 with the World Bank control and 0.891 with DESTA. The World Bank control
absorbs nearly the same variation as the regressor of interest, leaving a wide standard error; DESTA
overlaps less and the standard error narrows. Nothing about the estimated effect changes between the
two --- only how much of it the design can resolve.

The permutation test settles what the bootstrap cannot, because it reassigns whole agreement
profiles rather than partialling out a control and therefore gives the same verdict under both depth
measures. On the baseline
it returns $p = 0.278$ for the dirty coefficient and $p = 0.597$ for green; under DESTA, 0.146 and
0.480. A coefficient that would arise a quarter of the time from randomly relabelling which of these
destinations signed environmental chapters is not evidence about environmental content.

A leave-one-out exercise sharpens the diagnosis, and not in the usual way. The point estimate is
remarkably stable, sitting between $-0.0097$ and $-0.0133$ across all 23 exclusions. What is fragile
is precision: dropping Australia moves the coefficient by 13 percent but nearly triples its standard
error, from 0.0030 to 0.0087; dropping Korea doubles it. These two destinations are not outliers
pulling the estimate --- they supply the variation that identifies it. And which of the two is
pivotal depends on the depth control: under DESTA it is Korea rather than Australia.

\subsection*{Between firms or within them?}

Aggregating the data to product--destination--year cells and re-estimating with the corresponding
fixed effects yields a dirty coefficient of $-0.0119$ against $-0.0044$ in the firm-level panel, a
factor of 2.7. The two specifications differ by exactly one thing --- whether the firm dimension is
in the absorbers --- so the gap measures it: $(0.0119-0.0044)/0.0119 = 0.63$. Roughly three-fifths of
the aggregate dirty association is a change in \emph{which firms} export, not reallocation inside a
given firm.

This reconciles the present result with recent work on Chinese customs data reporting that
environmental provisions reduce dirty exports and raise green ones. That literature identifies from
variation across firms and destinations; the specification here removes it. The two are not in
conflict once the estimands are stated.

\begin{table}[H]
\centering
\footnotesize
\caption{The result in one table (log export value, World Bank depth control)}
\label{tab:summary}
\begin{threeparttable}
\begin{tabular}{lcccc}
\toprule
 & \multicolumn{2}{c}{WB EP depth} & \multicolumn{2}{c}{TREND EP count} \\
\cmidrule(lr){2-3}\cmidrule(lr){4-5}
Sample / design & $\times$ green & $\times$ dirty & $\times$ green & $\times$ dirty \\
\midrule
Full panel (reference)      & $-$0.0023 & $-$0.0044 & $-$0.0001 & $-$0.0009 \\
\quad bootstrap $p$         & [0.686] & [0.185] & [0.931] & [0.177] \\
\addlinespace[3pt]
Collapsed panel             & $-$0.0046 & $-$0.0119 & $+$0.0018 & $+$0.0004 \\
\quad bootstrap $p$         & [0.649] & [0.072] & [0.390] & [0.857] \\
\quad permutation $p$       & [0.597] & [0.278] & [0.160] & [0.817] \\
\addlinespace[3pt]
C-prod-HS4                  & $-$0.0009 & $+$0.0066 & $+$0.0005 & $+$0.0024 \\
\quad bootstrap $p$         & [0.860] & [0.435] & [0.730] & [0.734] \\
\addlinespace[3pt]
CEM-matched destinations    & $-$0.0023 & $-$0.0041 & $-$0.0006 & $-$0.0011 \\
\quad bootstrap $p$         & [0.610] & [0.251] & [0.609] & [0.204] \\
\addlinespace[3pt]
Deep vs.\ shallow (partners only) & $-$0.0022 & $-$0.0034 & $-$0.0004 & $-$0.0006 \\
\quad bootstrap $p$         & [0.585] & [0.499] & [0.751] & [0.739] \\
\bottomrule
\end{tabular}
\begin{tablenotes}[flushleft]\footnotesize
\item \emph{Notes:} Each cell is a separate estimation. Hong Kong and Macao excluded throughout; all
specifications control for non-environmental World Bank depth interacted with both product types.
Wild cluster bootstrap, $B = 9{,}999$, clustered by destination. Permutation $p$-values reported for
the collapsed panel, where the treated-only design is feasible. No stars: with this many designs,
star-counting invites exactly the selective reading the table is meant to prevent. The equivalent
table under the DESTA depth control is in the paper; the only cells that move are the dirty ones.
\end{tablenotes}
\end{threeparttable}
\end{table}


\subsection*{Does the content of the clauses matter?}

The natural objection to any null on an aggregate index is dilution. The indices add up every
environmental provision regardless of what it does, and most of what they add up is cooperation
language with no channel to a product mix.

The dilution is severe. Of 150 World Bank provision slots checked across the treated destinations,
8 --- 5.3 percent --- fall in sub-indices with a trade mechanism. In TREND it is 4 of 437, or 0.9
percent. A binding environmental obligation appears in exactly one of the fourteen agreements, the
2015 agreement with Korea. The two World Bank sub-indices with a trade mechanism are non-zero in only
three destination--years and always in the same proportion, which makes them perfectly collinear and
impossible to separate.

Replacing the aggregate index with one sub-index at a time reproduces the aggregate pattern: nothing
on the green margin anywhere, and on the dirty margin a marginally negative coefficient on
enforcement provisions that rests on asymptotic inference alone.

The conclusion is negative and it is the honest one. This sample cannot distinguish between ``there
is no effect'' and ``there are no provisions capable of producing one''. What it can establish is
that the two are observationally equivalent here.

\subsection*{Dynamics and the extensive margin}

Pre-treatment coefficients in the event study are flat on both margins. A mild negative drift for
green products late in the post-period is identified only by the cohorts that entered by 2010,
dominated by the two early ones (Bangkok 2002 and ASEAN 2005), and disappears
under the Sun--Abraham interaction-weighted estimator, whose aggregated effects are $-0.042$
($p = 0.27$) on green and $+0.073$ ($p = 0.28$) on dirty. Allowing destination-specific trends over
the full sample period turns the TREND green interaction significantly negative --- the only
green coefficient among the robustness checks on destination-specific trends to survive the bootstrap --- but the green gap was
already rising before entry into force, which makes a pre-existing trend the more economical account.

Poisson estimation on a zero-filled product--destination--year grid finds no green trade creation at
the extensive margin: $+0.0015$ ($p = 0.74$) with the World Bank index and $+0.0001$ ($p = 0.95$)
with TREND.

% ---------------------------------------------------------------------------
\section{Interpretation}

The result is not that environmental provisions cannot affect trade. It is that these provisions, in
these agreements, did not.

The literature that finds effects finds them in the same place. Brandi et al.\ (2020) report that
only provisions with a direct trade mechanism move anything: liberal provisions raise green export
shares, restrictive ones reduce dirty exports, and general cooperation language does nothing. Abman
et al.\ (2024) find that forest-preservation clauses offset the deforestation that typically follows
an agreement --- again, a specific clause with an identifiable mechanism. China's agreements of this
period are, on the same coding, overwhelmingly cooperation-only.

So the null is the predicted outcome of the existing literature's own heterogeneity findings, applied
to a sample where the treatment lacks the specificity that drives effects elsewhere. The policy
reading is direct: a chapter in an agreement is not the instrument. What the chapter contains is.

There is also a methodological contribution, arrived at reluctantly. With treatment concentrated
across roughly a dozen agreements, an asymptotically overwhelming composition effect ($p < 0.001$)
can dissolve entirely under bootstrap and permutation inference. Applying those tests as the primary
mode of inference rather than as a robustness appendix should be standard for this class of design.

% ---------------------------------------------------------------------------
\section{What the paper cannot say}

Four limitations, stated plainly.

\textbf{Levels.} The design speaks only to composition. Whether agreements with more environmental
content changed aggregate trade with those partners is not identified here, and cannot be: a
destination--year intercept absorbs the question along with the variable.

\textbf{The dirty margin is genuinely unresolved.} The coefficient is negative almost everywhere and
its statistical standing depends on a modelling choice. We read it as indistinguishable from a
placebo on the strength of the permutation test, but a dose--response interpretation --- the effect is
real and concentrated in the two destinations with the deepest environmental coverage --- is not ruled
out by the evidence assembled here.

\textbf{Provision content.} With thirteen distinct agreement profiles and perfect collinearity among
the trade-mechanism sub-indices, the claim that content matters is an inference across studies rather
than something this design demonstrates.

\textbf{Tariffs and firm entry.} The preferential tariff rate was not obtainable from the customs
data; the MFN rate serves as a control in the robustness battery only, and the headline estimates do
not use it. And the extensive margin is probed at the product--destination level: a new firm
beginning to export a green product that other Chinese firms already ship to that destination is not
detected.

\end{document}
